% =====================================================================
%  Where Is the Volatility Risk Premium?
%  Commodities, rates, credit and the limits of volatility timing
%
%  Compile from this folder:  latexmk -pdf main.tex
%  Numbers, tables and figures come from generated/ , written by
%      python src/run_all.py --data FILE            (smile strike)
%      python src/run_all.py --data FILE --strike atm
%      python src/summarise.py && python src/paper_exhibits.py
% =====================================================================
\documentclass[11pt,a4paper]{article}
\usepackage[margin=2.5cm]{geometry}
\usepackage[T1]{fontenc}
\usepackage{mathptmx}
\usepackage{amsmath,amssymb,amsthm}
\usepackage{booktabs}
\usepackage{graphicx}
\usepackage{caption}
\usepackage{xcolor}
\usepackage{enumitem}
\usepackage{setspace}
\usepackage{float}
\usepackage[authoryear,round]{natbib}
\usepackage[colorlinks=true,linkcolor=blue!50!black,citecolor=blue!50!black,urlcolor=blue!50!black]{hyperref}

\captionsetup{font=small,labelfont=bf,skip=4pt}
\setlist{itemsep=2pt,topsep=4pt}
\onehalfspacing

\newtheorem{hypothesis}{Hypothesis}
\newtheorem{proposition}{Proposition}
\theoremstyle{remark}
\newtheorem*{remark}{Remark}

\newcommand{\E}{\mathbb{E}}
\newcommand{\Var}{\mathrm{Var}}
\newcommand{\F}{\mathcal{F}}
\newcommand{\note}[1]{\par\smallskip\parbox{\linewidth}{\footnotesize #1}}

\input{generated/numbers}

\title{Where Is the Volatility Risk Premium?\\[0.3em]
\large Commodities, Rates, Credit and the Limits of Volatility Timing}
\author{Paolo Mastrogiacomo\thanks{EDHEC Business School. Supervisor: Prof.\ Enrique Schroth. Preliminary draft; comments welcome. All numbers are produced by the code in the accompanying repository.}}
\date{\today}

\begin{document}
\maketitle

\begin{abstract}
\noindent
I sell one month of variance on \NInstr{} instruments in five asset classes (energy, metals, grains, government bonds and credit) from 2007 to 2026, with one strategy, one set of timing rules and one inference framework, and use the S\&P~500 as a benchmark. The variance risk premium is not a general property of option markets. It is large and statistically robust in bond futures (Sharpe ratios $\SRThirtyy$ to $\SRBund$, $t$-statistics of the volatility premium between \VPtThirtyy{} and \VPtBund), in the CDX investment-grade index ($t=\VPtCdx$), and in corn and soybeans ($t=\VPtCorn$ and \VPtSoy), comparable to the equity premium (Sharpe $\SRSpx$). It is absent in crude oil, natural gas, gold, silver, copper, wheat and the two credit ETFs. Equity-style timing rules do not improve the strategy outside equities and credit: \NSig{} of \NTests{} changes in Sharpe ratio are significant at 5\%, against \NExpected{} expected by chance, and none survives a Holm correction. A mean--variance framework explains where timing can work: de-risking after volatility rises pays only when the conditional Sharpe ratio falls with volatility. That is the case for the S\&P~500 and CDX IG, and in no other market in the sample.
\end{abstract}

\bigskip
\noindent\textbf{Keywords:} variance risk premium, volatility timing, commodity futures, bond futures, credit indices, volatility-managed portfolios.\\
\noindent\textbf{JEL:} G11, G12, G13, Q02.

\newpage
\section{Introduction}

Option-implied variance has historically exceeded the variance subsequently realised, so investors who sell variance earn a premium \citep{carrwu2009,bollerslev2009}. The return stream is asymmetric: frequent small gains and rare, very large losses \citep{duarte2007}. Two questions follow for anyone who harvests it. Where does the premium exist? And can its losses be avoided by reducing exposure when volatility rises, as the volatility-management literature suggests \citep{moreiramuir2017,yang2024}?

Most of what we know comes from equity indices. The premium has been documented in commodities \citep{trolle2010,prokopczuk2017}, in interest-rate options \citep{trolleschwartz2014,choi2017} and across asset classes \citep{fallon2015,khalfallah2025}, but with different data, strikes, sample periods and test statistics, which makes comparisons across markets hard. Volatility timing has been tested almost exclusively on equity factors and equity volatility, and its gains are fragile out of sample \citep{cederburg2020}.

This paper applies one design to all markets. At each month-end I sell a synthetic one-month variance swap on \NInstr{} instruments: eight commodity futures, five government bond futures, two credit ETFs and the CDX North America investment-grade index. The S\&P~500 serves as the benchmark. The strike is the implied variance integrated over the option smile where the smile is quoted close enough to the money; otherwise it is at-the-money implied variance, a lower bound. For CDX IG and the S\&P~500 it is the model-free Credit VIX and VIX. I then apply seven pre-specified timing rules and test every rule in every market.

The results are summarised in Figure~\ref{fig:premium}. Three findings stand out.

\begin{enumerate}
\item \textbf{The premium is concentrated.} It is large and significant in all five bond futures, in CDX IG and in corn and soybeans, with Sharpe ratios of $\SRThirtyy$ to $\SRCorn$ that bracket the $\SRSpx$ of the S\&P~500 under the same method. It is insignificant in energy, metals, wheat and the credit ETFs, whose strategies lose money in the crises of 2008 and 2020.
\item \textbf{Equity-style timing does not travel.} Of \NTests{} rule-by-market changes in Sharpe ratio, \NSig{} are significant at 5\%, \NSigUp{} improvements and \NSigDown{} deteriorations, close to the \NExpected{} expected under the null. The exceptions are CDX IG and, less strongly, the S\&P~500.
\item \textbf{The framework predicts where timing works.} Proposition~\ref{prop:weight} shows that the case for de-risking after volatility rises depends on how the premium scales with variance. Empirically, the conditional Sharpe ratio falls with volatility only in CDX IG and the S\&P~500, exactly the two markets where timing helps. In bonds and grains it is flat or rising, and timing adds nothing.
\end{enumerate}

The contribution is a like-for-like map of the variance risk premium across asset classes and a test of volatility timing that is disciplined by a mechanism. The paper also shows how much of the premium survives realistic costs, and that its cross-asset pattern is consistent with a hedging-demand explanation: the premium exists where options are used by natural hedgers \citep{bollenwhaley2004,garleanu2009}. That explanation is formulated after seeing the data, and Section~\ref{sec:next} sets out how to test it out of sample.

Section~\ref{sec:lit} reviews the literature. Section~\ref{sec:theory} sets out the framework. Sections~\ref{sec:data} and~\ref{sec:method} describe the data and the method. Section~\ref{sec:premium} reports the premium, Section~\ref{sec:timing} the timing results, Section~\ref{sec:mech} the mechanism and Section~\ref{sec:portfolio} diversification. Section~\ref{sec:discussion} discusses the results and Section~\ref{sec:next} the limitations and next steps.

\section{Related literature}
\label{sec:lit}

\paragraph{Variance risk premia.} The variance risk premium is the difference between the risk-neutral and the physical expectation of future variance. \citet{carrwu2009} measure it with synthetic variance swaps and \citet{bollerslev2009} show that it predicts equity returns. \citet{bakshikapadia2003} document negative returns on delta-hedged options. The model-free implied variance follows \citet{brittenjones2000}, \citet{demeterfi1999} and \citet{jiangtian2005}.

\paragraph{Beyond equities.} \citet{trolle2010} document variance risk premia in energy and \citet{prokopczuk2017} across commodities. \citet{trolleschwartz2014} and \citet{choi2017} find variance risk premia in interest-rate swaptions and Treasury futures options. \citet{fallon2015} show that volatility premia exist across asset classes and diversify, and \citet{khalfallah2025} study volatility carry across asset classes. This paper uses a common design to compare these markets directly, adds credit through the CDX index and the Credit VIX \citep{spdji_creditvix}, and reports where the premium is absent as well as where it is present.

\paragraph{Why option prices carry a premium.} In demand-based option pricing, end users' net demand for options raises their prices when dealers cannot hedge perfectly \citep{bollenwhaley2004,garleanu2009}. In futures markets, hedging pressure explains part of the risk premium \citep{deroon2000}, and the capital of intermediaries prices risk across asset classes \citep{hekellymanela2017}. These papers motivate the interpretation in Section~\ref{sec:discussion}.

\paragraph{Volatility management.} \citet{moreiramuir2017} and \citet{barroso2015} scale exposure by inverse variance. \citet{yang2024} applies the idea to volatility assets and attributes the gains to a conditional premium that does not keep up with volatility. \citet{cederburg2020} show that much of the gain disappears out of sample. Testing many rules raises the risk of false discoveries \citep{white2000,hansen2005,romanowolf2005,harvey2016}.

\section{Economic framework}
\label{sec:theory}

Let $RV_{t+1}$ be the variance realised over the month after date $t$ and $K_t$ the variance-swap strike, equal to the risk-neutral expectation $\E^{\mathbb{Q}}_t[RV_{t+1}]$. A seller of one unit of variance notional earns
\begin{equation}
x_{t+1} = K_t - RV_{t+1},
\qquad
\mu_t \equiv \E_t[x_{t+1}] = \E^{\mathbb{Q}}_t[RV_{t+1}] - \E^{\mathbb{P}}_t[RV_{t+1}].
\label{eq:x}
\end{equation}
Let $v_t = \E^{\mathbb{P}}_t[RV_{t+1}]$ be the conditional variance and assume
\begin{align}
\text{(A1)}\quad & \Var_t(x_{t+1}) = \kappa\, v_t^{2}, \label{eq:a1}\\
\text{(A2)}\quad & \mu_t = \lambda\, v_t^{\,b}, \qquad \lambda>0. \label{eq:a2}
\end{align}
(A1) states that the volatility of variance is proportional to its level \citep{andersen2003}; (A2) lets the premium scale with the state, with elasticity $b$.

\begin{proposition}
\label{prop:weight}
For a mean--variance investor with risk aversion $\gamma$, the optimal exposure to short variance is $w_t^{*}=\mu_t/(\gamma\Var_t(x_{t+1}))$. Under (A1) and (A2),
\begin{equation}
w_t^{*} = \frac{\lambda}{\gamma\kappa}\; v_t^{\,b-2},
\qquad
SR_t = \frac{\lambda}{\sqrt{\kappa}}\; v_t^{\,b-1}.
\label{eq:wstar}
\end{equation}
The optimal exposure falls when volatility rises if and only if $b<2$, and the conditional Sharpe ratio falls when volatility rises if and only if $b<1$.
\end{proposition}

The practitioner view (``sell after spikes'') corresponds to $b>2$; constant exposure is optimal at $b=2$; inverse-variance scaling \citep{moreiramuir2017} is optimal at $b=1$, where the conditional Sharpe ratio is constant. The lower $b$, the more the premium lags behind risk and the larger the gain from cutting exposure after volatility rises. When $b<1$ the conditional Sharpe ratio itself falls in volatile months, and de-risking removes months that are both risky and poorly paid. Dividing (\ref{eq:a2}) by $v_t$ gives a testable form:
\begin{equation}
\frac{x_{t+1}}{v_t} = \lambda\, v_t^{\,b-1} + \varepsilon_{t+1},
\label{eq:test}
\end{equation}
whose residual is homoskedastic under (A1). A negative relation between $x_{t+1}/v_t$ and $\ln v_t$ indicates $b<1$.

\subsection*{Hypotheses}

\begin{hypothesis}[Premium]\label{h:premium}
In each market, implied variance exceeds subsequent realised variance on average: $\E[x_{t+1}]>0$.
\end{hypothesis}
\begin{hypothesis}[Tails]\label{h:tail}
Losses of the short-variance position are concentrated in a few months of market-wide stress.
\end{hypothesis}
\begin{hypothesis}[Timing]\label{h:timing}
Rules that reduce exposure after volatility rises improve risk-adjusted performance.
\end{hypothesis}
\begin{hypothesis}[Mechanism]\label{h:mech}
Timing gains are larger where the conditional Sharpe ratio falls with volatility ($b<1$).
\end{hypothesis}
\begin{hypothesis}[Diversification]\label{h:div}
A portfolio of short-variance positions across markets has a smaller tail than its components.
\end{hypothesis}

The rules, thresholds and hypotheses \ref{h:premium}--\ref{h:timing} were fixed in a research design written for commodities before any result was computed; the same design was then applied unchanged to rates, credit and the benchmark.

\section{Data}
\label{sec:data}

\subsection{Instruments and sample}

Daily Bloomberg data run from January 2007 to October 2026 (Table~\ref{tab:sample}). For each futures market I use the first and second generic contracts and the 30-day implied volatility at 100\% moneyness; for commodities also at 90\% and 110\%. For the credit ETFs I use the price and the same implied-volatility fields. For CDX IG I use the on-the-run 5-year spread and the CDX/Cboe IG Credit VIX (VIXIG), which is expressed in basis points of spread volatility and starts in 2012 \citep{spdji_creditvix}. For the S\&P~500 I use the index and the VIX.

\begin{table}[H]
\centering
\caption{Instruments and sample}
\label{tab:sample}
\small
\input{generated/tab_sample}
\note{Months: decision months with a complete following month. Mean IV and RV: average implied volatility at decision dates and realised volatility over the following month, in \% per year (CDX IG in basis points of spread). Strike: variance-swap strike (Section~\ref{sec:iv}).}
\end{table}

\subsection{Returns}

\paragraph{Futures.} Generic futures switch contract at each expiry, so I construct the return of a position that always holds a single contract. For each contract I rebuild the last trading day and, for physically delivered contracts, the first notice day from exchange rules: NYMEX, COMEX and CBOT for commodities, CBOT for Treasury futures (last business day of the delivery month for the 2- and 5-year notes; seventh business day before it for the 10-year note and the bond) and Eurex for the Bund (two exchange days before delivery on the 10th). The position switches from $F_1$ to $F_2$ five trading days before the earlier of the two dates. The calendar is validated on the data: the change in $\ln(F_1/F_2)$ spikes on the day after each rebuilt last trading day for every contract (Appendix Table~\ref{tab:calendar}).

\paragraph{ETFs and indices.} HYG, LQD and the S\&P~500 use daily log price returns. For CDX IG, the return is the daily change in the spread in basis points, so that realised volatility is in the same units as VIXIG. The on-the-run index switches series on the first trading day on or after 20 March and 20 September; those days are excluded.

\paragraph{Realised variance.} $RV_{t+1} = \frac{252}{n}\sum_{i=1}^{n} r_i^2$ over the $n$ daily returns between consecutive month-ends.

\subsection{Implied variance}
\label{sec:iv}

The fair variance-swap strike is the model-free implied variance \citep{brittenjones2000,demeterfi1999}. For commodities I integrate undiscounted Black prices of out-of-the-money options over the 30-day smile, interpolating volatility linearly between 90\%, 100\% and 110\% moneyness and holding it flat beyond \citep{jiangtian2005}. With three points this captures the slope and curvature of the smile near the money but not its far tails, so it remains a lower bound. For bond futures and ETFs, 90\% and 110\% moneyness lie many standard deviations from the money and I use ATM implied variance, also a lower bound. VIXIG and VIX are themselves model-free. Table~\ref{tab:robust} reports commodity results with the ATM strike as well.

\section{Method}
\label{sec:method}

\paragraph{Strategy.} At each month-end $t$ the investor sells one month of variance at strike $K_t$. Per unit of vega notional the P\&L is $(K_t-RV_{t+1})/(2\sqrt{K_t})$, so one volatility point of premium is worth about 1\% of notional. All strategies are scaled to 10\% annual volatility for comparison.

\paragraph{Timing rules.} Seven rules map information available at $t$ into a weight $w_t$ (Table~\ref{tab:rules}). Rules that need a history use expanding windows and keep $w_t=1$ for the first 12 months.

\begin{table}[H]
\centering
\caption{Timing rules}
\label{tab:rules}
\small
\begin{tabular}{lll}
\toprule
Rule & Signal at $t$ & Weight \\
\midrule
Volatility scaling & Past-month realised variance $RV_t$ & $\min\{2,\ \tilde v_t / RV_t\}$ \\
Realised-vol threshold & $RV_t$ above its 80th percentile & 0.5 \\
Implied-vol jump & Implied vol up more than 20\% in 5 days & 0 \\
Implied-vol trend & Implied vol above its 3-month average & 0.5 \\
Term structure & 1st/2nd-month implied vol above its 80th percentile & 0.5 \\
Premium filter & Implied vol below past realised vol & 0 \\
Backwardation & $\ln(F_1/F_2)$ above its 80th percentile & 0.5 \\
\bottomrule
\end{tabular}
\note{$\tilde v_t$: expanding median of past realised variance. Otherwise $w_t=1$. The term-structure rule needs a second implied-vol maturity (not available for VIXIG and VIX) and the backwardation rule a futures curve (not available for ETFs and indices).}
\end{table}

\paragraph{Inference.} Means are tested with \citet{neweywest1987} standard errors (3 lags). Each managed strategy is compared with the unmanaged one at equal volatility with a spanning regression and with a circular block bootstrap of the difference in Sharpe ratios \citep{politis1992,ledoitwolf2008}, with 6-month blocks and 5{,}000 draws. Because \NTests{} rule-by-market tests are run, I report how many are significant against the number expected by chance, and the \citet{holm1979} step-down correction.

\paragraph{Costs.} A fixed cost in volatility points is not comparable between a 1.4\% bond future and 50\% natural gas. Costs are therefore charged as a fraction of implied volatility per monthly trade: 2.5\%, 5\% and 10\%.

\section{The premium across asset classes}
\label{sec:premium}

\begin{figure}[H]
\centering
\includegraphics[width=\linewidth]{generated/fig_premium.pdf}
\caption{Volatility premium and Sharpe ratio of short one-month variance, 2007--2026. Left: implied minus subsequently realised volatility in \% of implied volatility, with 95\% Newey--West confidence intervals. Right: annualised Sharpe ratio; dashed line: S\&P~500 with the VIX strike. CDX IG from 2012.}
\label{fig:premium}
\end{figure}

Table~\ref{tab:premium} reports the premium. Hypothesis~\ref{h:premium} holds in eight of sixteen markets.

\begin{itemize}
\item \textbf{Bond futures.} All five carry a significant premium, between \RelThirtyy\% and \RelTwoy\% of implied volatility, with $t$-statistics from \VPtThirtyy{} to \VPtTwoy{} and Sharpe ratios from $\SRThirtyy$ to $\SRBund$. The result holds in the United States and Germany and with a strike that understates the premium.
\item \textbf{CDX IG.} Implied volatility exceeds realised by \VPCdx{} basis points a year, \RelCdx\% of implied volatility ($t=\VPtCdx$), and the premium is positive in \PosCdx\% of months and in all \CdxYears{} calendar years (between \CdxYearMin{} and \CdxYearMax{} basis points).
\item \textbf{Grains.} Corn and soybeans earn \RelCorn\% and \RelSoy\% of implied volatility ($t=\VPtCorn$ and \VPtSoy), with the highest Sharpe ratios in the sample ($\SRCorn$ and $\SRSoy$) and the mildest skewness. Wheat does not ($t=\VPtWheat$).
\item \textbf{No premium.} In crude oil, natural gas, gold, silver and copper the premium is below 5\% of implied volatility and insignificant, and the strategies have negative Sharpe ratios. The credit ETFs show none either: HYG is insignificant ($t=\VPtHyg$) and LQD is zero.
\end{itemize}

The equity benchmark earns \RelSpx\% of implied volatility ($t=\VPtSpx$) and a Sharpe ratio of $\SRSpx$. The premium in bonds, CDX IG, corn and soybeans is therefore of the same order of magnitude as in equities.

\begin{table}[H]
\centering
\caption{Volatility premium and performance of short one-month variance}
\label{tab:premium}
\small
\input{generated/tab_premium}
\note{Vol premium: mean of implied minus realised volatility, in volatility points (CDX IG: basis points). $t$: Newey--West, 3 lags; ``variance'': the same test on the variance premium $K_t-RV_{t+1}$. Sharpe ratio annualised; skewness, maximum drawdown and 99\% CVaR of monthly returns scaled to 10\% annual volatility.}
\end{table}

\subsection{Tails, crises and seasonality}

Hypothesis~\ref{h:tail} holds everywhere (Table~\ref{tab:crises}, Figure~\ref{fig:cum}). Losses cluster in 2008 and 2020: COVID cost the crude oil strategy \CovidOil\% of capital at 10\% volatility, LQD \CovidLqd\% and CDX IG \CovidCdx\%. For CDX IG, March 2020 alone lost as much as \CdxWorstMultiple{} months of average premium. Markets with a premium differ from those without not in whether losses come, but in whether the premium earned between crises pays for them: the bond and grain strategies recover their losses within one to two years, while energy and metals give back years of gains and do not recover them.

\begin{figure}[H]
\centering
\includegraphics[width=\linewidth]{generated/fig_cumulative.pdf}
\caption{Cumulative P\&L of short one-month variance, scaled to 10\% annual volatility. Asset-class lines are equal-weighted averages of their instruments; credit ETFs are HYG and LQD.}
\label{fig:cum}
\end{figure}

\begin{table}[H]
\centering
\caption{Cumulative P\&L in stress episodes (\% of capital, 10\% annual volatility)}
\label{tab:crises}
\small
\input{generated/tab_crises}
\note{Episodes by month of realisation: GFC Sep--Dec 2008; euro debt crisis Aug--Nov 2011; taper tantrum May--Jun 2013; commodity crash Jul 2014--Feb 2016; COVID Feb--Apr 2020; Ukraine Feb--Jun 2022; tariffs Mar--Apr 2025. --: before the start of the series.}
\end{table}

The grain premium is seasonal (Figure~\ref{fig:seasonal}). It averages \WinCorn{} volatility points for corn from November to February and turns negative in March ($\MarCorn$) and June ($\JunCorn$), the months of the USDA Prospective Plantings, Acreage and Grain Stocks reports. Realised volatility then exceeds what options price, which suggests that options underprice scheduled supply information.

\begin{figure}[H]
\centering
\includegraphics[width=\linewidth]{generated/fig_seasonal.pdf}
\caption{Average volatility premium by month of realisation, grains. Shaded: March and June, the months of the main USDA reports.}
\label{fig:seasonal}
\end{figure}

\subsection{Robustness and costs}

Table~\ref{tab:robust} varies the specification. The Sharpe ratio in bonds, CDX IG, corn and soybeans is positive in both halves of the sample, does not depend on the roll date, and survives excluding March--April 2020; for CDX IG the Sharpe ratio rises to $\SRexcCdx$ without those two months. With the ATM strike, the commodity premium falls (corn $\SRatmCorn$, soybeans $\SRatmSoy$): part of it is the price of the smile, which is why the strike matters for any comparison. The performance of CDX IG and most bond futures weakens after 2016, while corn strengthens.

Costs are the binding constraint. At 2.5\% of implied volatility per trade, the bond strategies keep Sharpe ratios of $\SRcostAThirtyy$ to $\SRcostABund$; at 5\%, only CDX IG ($\SRcostBCdx$), corn ($\SRcostBCorn$) and soybeans ($\SRcostBSoy$) remain positive. CDX IG is the only one of the \NInstr{} markets that survives a 10\% cost ($\SRcostCCdx$).

\begin{table}[H]
\centering
\caption{Robustness: Sharpe ratio of short one-month variance}
\label{tab:robust}
\small
\input{generated/tab_robust}
\note{Baseline: smile strike for commodities, ATM for bonds and ETFs, VIXIG and VIX. ATM strike: shown where the baseline uses the smile. Roll $k$: switch to the second contract $k$ trading days before the roll date (baseline 5). Costs: per monthly trade, in \% of implied volatility.}
\end{table}

\section{Does volatility timing help?}
\label{sec:timing}

Figure~\ref{fig:timing} reports the change in Sharpe ratio from each rule in each market, and Table~\ref{tab:timing} summarises it by rule. Hypothesis~\ref{h:timing} is rejected outside equities and credit. Of \NTests{} tests, \NSig{} are significant at 5\%, \NSigUp{} improvements and \NSigDown{} deteriorations; under the null about \NExpected{} would be significant (binomial $p=\BinomP$). The smallest $p$-value is \MinP{}, and no test survives the \citet{holm1979} correction. Average changes are small: inverse-variance scaling, the central rule of the literature, raises the Sharpe ratio in \Upvolscaling{} of \Nvolscaling{} markets and lowers or leaves it unchanged in the rest.

The exceptions are informative. In CDX IG four rules improve the Sharpe ratio significantly; the implied-vol jump rule raises it from $\SRCdx$ to $\SRCdxIVJ$, because VIXIG more than doubled in the last week of February 2020, which triggered the exit just before the March 2020 losses. In the S\&P~500, volatility scaling raises the Sharpe ratio from $\SRSpx$ to $\SRSpxVS$ ($p=\PSpxVS$), in line with \citet{yang2024}. Section~\ref{sec:mech} shows why these two markets differ from the others.

\begin{figure}[H]
\centering
\includegraphics[width=\linewidth]{generated/fig_timing.pdf}
\caption{Change in Sharpe ratio from each timing rule, relative to the unmanaged strategy at equal volatility. * and bold: significant at 5\% (circular block bootstrap). --: rule not applicable.}
\label{fig:timing}
\end{figure}

\begin{table}[H]
\centering
\caption{Timing rules: summary across markets}
\label{tab:timing}
\small
\input{generated/tab_timing}
\note{$\Delta SR$: change in annualised Sharpe ratio relative to the unmanaged strategy. Better/Worse: number of markets in which the change is significant at 5\%.}
\end{table}

\section{Mechanism: when does de-risking pay?}
\label{sec:mech}

Proposition~\ref{prop:weight} says that de-risking after volatility rises pays when the premium lags behind risk, most clearly when the conditional Sharpe ratio falls with volatility ($b<1$). Table~\ref{tab:mechanism} tests this with equation~(\ref{eq:test}), proxying $v_t$ by past-month realised variance, and with conditional Sharpe ratios by quintile of past variance (Figure~\ref{fig:quintiles}).

The slope of $x_{t+1}/v_t$ on $\ln v_t$ is significantly negative in exactly two markets: CDX IG ($t=\SlopeTCdx$) and the S\&P~500 ($t=\SlopeTSpx$). In CDX IG the conditional Sharpe ratio falls from $\QoneSRCdx$ in the calmest quintile to $\QfiveSRCdx$ in the most volatile; in the S\&P~500 from $\QoneSRSpx$ to $\QfiveSRSpx$. These are the two markets where timing rules help. In bonds and grains the conditional Sharpe ratio is flat or rises with volatility (US 30-year: $\QoneSRThirtyy$ to $\QfiveSRThirtyy$; soybeans: $\QoneSRSoy$ to $\QfiveSRSoy$), $b=1$ cannot be rejected, and timing adds nothing. Hypothesis~\ref{h:mech} is supported: the mechanism estimated without any timing rule predicts where the rules work.

The proxy for $v_t$ is noisy, which biases the slope towards zero; the tests therefore have little power to detect moderate departures from $b=1$. The quintile sorts, which use only the ranking of past variance, give the same picture.

\begin{table}[H]
\centering
\caption{How the premium scales with volatility}
\label{tab:mechanism}
\small
\input{generated/tab_mechanism}
\note{Regression of $x_{t+1}/v_t$ on $\ln v_t$, with $x_{t+1}=K_t-RV_{t+1}$ and $v_t$ the realised variance of the past month; Newey--West $t$-statistics. A negative slope implies $b<1$ (conditional Sharpe ratio falls with volatility). SR: annualised Sharpe ratio of $x_{t+1}$ within quintiles of $v_t$.}
\end{table}

\begin{figure}[H]
\centering
\includegraphics[width=\linewidth]{generated/fig_quintiles.pdf}
\caption{Conditional Sharpe ratio of short variance by quintile of past-month realised variance. Thin lines: instruments; thick lines: asset-class average.}
\label{fig:quintiles}
\end{figure}

\section{Diversification}
\label{sec:portfolio}

Short-variance strategies are only moderately correlated across markets: the average pairwise correlation of monthly returns is \AvgCorr{}. Grains are almost uncorrelated with bonds (\CorrGrainsRates) and credit (\CorrGrainsCredit), while bonds and credit move together (\CorrRatesCredit) and with equities (\CorrRatesEq{} and \CorrCreditEq), because they lose in the same episodes of market stress.

An equal-weighted portfolio of all \NInstr{} strategies earns a Sharpe ratio of $\SRAll$ (Table~\ref{tab:portfolio}). It includes the markets without a premium, so it is the relevant benchmark for an investor who has no prior on where the premium is. Its performance falls from $\SRAllA$ in 2007--2016 to $\SRAllB$ in 2017--2026, and its skewness remains $\SkewAll$: diversification smooths the path but does not remove the common tail, so Hypothesis~\ref{h:div} holds only partly.

At the portfolio level the implied-vol jump rule raises the Sharpe ratio from $\SRAll$ to $\SRAllIVJ$ ($p=\PAllIVJ$), in both halves of the sample. Most of the gain comes from CDX IG: excluding it, the rule raises the Sharpe ratio from $\SRExCdx$ to $\SRExCdxIVJ$, which is not significant ($p=\PExCdxIVJ$). Volatility scaling does not help ($\SRAllVS$, $p=\PAllVS$).

\begin{table}[H]
\centering
\caption{Equal-weighted portfolios of short-variance strategies}
\label{tab:portfolio}
\small
\input{generated/tab_portfolio}
\note{Equal weights on the strategies available in each month (CDX IG from 2012), each scaled to 10\% annual volatility; portfolios rescaled to 10\%.}
\end{table}

\section{Discussion}
\label{sec:discussion}

The markets with a premium share one feature: options are bought by natural hedgers who need the protection whatever its price. Banks, insurers and pension funds hedge interest-rate and credit risk with bond-futures options and credit index options; grain producers and processors hedge the crop with options around planting and harvest. In these markets, dealers who sell the options are paid a premium for bearing variance risk they cannot hedge \citep{bollenwhaley2004,garleanu2009}. In energy and metals, option demand is more balanced between producers, consumers and speculators, and volatility is dominated by supply shocks that options appear to price fully. The credit ETFs are a useful contrast: their options trade mostly among investors, and their realised volatility includes the liquidity crises of the ETFs themselves.

This interpretation is formed after seeing the data and is not yet tested. It makes predictions that can be tested out of sample (Section~\ref{sec:next}): the premium should be larger where hedgers' net option demand is larger \citep{deroon2000} and when intermediaries have less capital \citep{hekellymanela2017}.

For investors, the results suggest that where to sell volatility matters more than when. A diversified short-volatility allocation should weight bonds, investment-grade credit and grains, not energy and metals. Timing rules built on equity experience should not be applied elsewhere without checking how the conditional premium scales with volatility; in credit, an exit after a sharp rise in implied volatility is the one rule with economic support.

\section{Limitations and next steps}
\label{sec:next}

\paragraph{Limitations.} The variance swap is synthetic: commodity and bond variance swaps are thinly traded and the strategy would be implemented with listed options, at costs that Table~\ref{tab:robust} shows to be decisive. The strike for bonds and ETFs is ATM and for commodities uses three smile points; both understate the true variance strike. The CDX result rests on one index since 2012, and the timing gains in credit on one episode. Rules are rebalanced monthly.

\paragraph{Next steps.}
\begin{enumerate}
\item \textbf{Model-free strike for bonds and ETFs:} 30-day implied volatility at 95\%, 97.5\%, 102.5\% and 105\% moneyness.
\item \textbf{More credit:} iTraxx Main (VIXIE), iTraxx Crossover (VIXXO), CDX HY (VIXHY) and the percentage version VIXIGP, so that the credit result does not rest on one index.
\item \textbf{Out-of-sample test of the hedging-demand interpretation:} before running the code, record a prediction for markets not yet used (sugar, coffee, cotton, live cattle, lean hogs: premium; Brent, heating oil: none), then test it.
\item \textbf{Mechanism:} relate the premium to hedgers' net option positions (CFTC Commitments of Traders, futures and options combined), option open interest and intermediary capital \citep{hekellymanela2017}; test the grain seasonality on USDA report dates.
\item \textbf{Implementation:} bid--ask quotes on one-month ATM options to measure costs; a delta-hedged straddle as a traded counterpart; a HAR forecast of $v_t$ and data-snooping tests on the full set of rules \citep{white2000,hansen2005,romanowolf2005}.
\end{enumerate}

\section*{Data and code availability}

Futures, option-implied volatility and index data are from Bloomberg and cannot be redistributed. All code, including the contract-calendar reconstruction, the strategies, the timing rules, the tests and the generation of every table, figure and number in this paper, is in the accompanying repository, with a synthetic-data generator that runs the pipeline without licensed data.

\bigskip
\noindent\textit{Declaration of AI use.} [Tool name, date] was used for [purpose]. \textit{(To be completed according to EDHEC guidelines and agreed with the supervisor.)}

\bibliographystyle{apalike}
\bibliography{references}

\appendix
\section{Additional tables}

\begin{table}[H]
\centering
\caption{Validation of the contract calendars}
\label{tab:calendar}
\small
\input{generated/tab_calendar}
\note{Median absolute daily change in $\ln(F_1/F_2)$ on days around each rebuilt last trading day, relative to the median over all days. Bloomberg generics switch contract on the day after the last trading day, so the ratio should peak at offset $+1$, as it does for every contract.}
\end{table}

\begin{table}[H]
\centering
\caption{Correlation of monthly returns of asset-class portfolios}
\label{tab:corr}
\small
\input{generated/tab_corr}
\note{Equal-weighted short-variance portfolios by asset class; Equity: S\&P~500 with the VIX strike.}
\end{table}

\end{document}
